% !TeX root = ./main.tex
% Acronyms and terms -- maintained BY HAND, sorted AUTOMATICALLY when printed.
% Type them as plain text in the body; nothing expands or links. Two entry
% types, rendered alike and sorted case-insensitively into one list, so the
% order here does not matter:
%
%   \acr{PET}{\textbf{positron emission tomography}---a functional modality}
%   \term{positron range}{the displacement between emission and annihilation}
%
% Sorted with expl3, not datatool: its load-and-sort cost helped push Overleaf
% past its compile timeout. Expandable string compare is \pdfstrcmp on pdfTeX
% (Overleaf), \strcmp on XeTeX (tectonic).
\ifdefined\pdfstrcmp\let\acrstrcmp\pdfstrcmp\else\let\acrstrcmp\strcmp\fi
\ExplSyntaxOn
\seq_new:N \g__acr_seq
\bool_new:N \g__acr_first_bool
% Case fold for the sort key. Renamed across expl3 releases, so take whichever
% this kernel has.
\cs_if_exist:NTF \str_casefold:n
	{ \cs_new_eq:NN \__acr_fold:n \str_casefold:n }
	{
	\cs_if_exist:NTF \str_foldcase:n
		{ \cs_new_eq:NN \__acr_fold:n \str_foldcase:n }
		{ \cs_new_eq:NN \__acr_fold:n \tl_to_str:n }
	}
\cs_new_protected:Npn \acr #1#2
	{ \seq_gput_right:Nn \g__acr_seq { {#1} {#2} } }
% Same mechanism; the separate name only records which kind an entry is.
\cs_new_eq:NN \term \acr
\cs_new_protected:Npn \__acr_row:nn #1#2
	{ \textbf{#1} & #2 }
\cs_new_protected:Npn \printacronyms
	{
	\seq_sort:Nn \g__acr_seq
		{
		\int_compare:nNnTF
			{ \acrstrcmp
				{ \exp_args:Ne \__acr_fold:n { \use_i:nn ##1 } }
				{ \exp_args:Ne \__acr_fold:n { \use_i:nn ##2 } } }
			> { 0 }
			{ \sort_return_swapped: }
			{ \sort_return_same: }
		}
	\bool_gset_true:N \g__acr_first_bool
	\seq_map_inline:Nn \g__acr_seq
		{
		\bool_if:NTF \g__acr_first_bool
			{ \bool_gset_false:N \g__acr_first_bool }
			{ \\ }
		\__acr_row:nn ##1
		}
	}
\ExplSyntaxOff

\acr{PET}{\textbf{positron emission tomography}---a functional imaging modality}
\acr{CT}{\textbf{computed tomography}---an anatomical imaging modality}
\acr{MRI}{\textbf{magnetic resonance imaging}---an anatomical imaging modality}
\acr{FDG}{\textbf{fluorodeoxyglucose}---a glucose-analogue radiotracer}
\acr{LOR}{\textbf{line of response}---the straight path along which a coincident pair of gamma photons are inferred to have originated. Can alternatively be used to refer to the straight path which connects two activated detectors of a PET scanner}
\acr{TOF}{\textbf{time-of-flight}---a technique using high-speed timing to better localise the annihilation position along a LOR}
\acr{SNR}{\textbf{signal-to-noise ratio}---a measure of data quality}
\acr{ICS}{\textbf{inter-crystal scatter}---the phenomenon where a gamma photon Compton scatters in one segment of a PET detector and goes on to interact inside another nearby segment}
\acr{LSO}{\textbf{lutetium oxyorthosilicate}---an inorganic scintillator crystal}
\acr{LYSO}{\textbf{lutetium-yttrium oxyorthosilicate}---an inorganic scintillator crystal}
\acr{BGO}{\textbf{bismuth germanate}---an inorganic scintillator crystal}
\acr{NaI(Tl)}{\textbf{thallium-doped sodium iodide}---an inorganic scintillator crystal}
\acr{GAGG}{\textbf{gadolinium aluminium gallium garnet}---an inorganic scintillator crystal}
\acr{CZT}{\textbf{cadmium zinc telluride}---a semiconductor detector material, converting gamma photons to charge directly rather than via scintillation}
\acr{PMT}{\textbf{photomultiplier tube}---a vacuum photodetector, the conventional readout for scintillators}
\acr{SiPM}{\textbf{silicon photomultiplier}---a solid-state photodetector, an array of SPADs read out in parallel}
\acr{SPAD}{\textbf{single-photon avalanche diode}---the individual sensing element of a SiPM}
\acr{APD}{\textbf{avalanche photodiode}---a solid-state photodetector operated below breakdown}
\acr{ASIC}{\textbf{application-specific integrated circuit}---custom readout electronics}
\acr{ADC}{\textbf{analogue-to-digital converter}}
\acr{TDC}{\textbf{time-to-digital converter}}
\acr{DAQ}{\textbf{data acquisition}}
\acr{CTR}{\textbf{coincidence resolving time}---the timing resolution of a detector pair, which sets the achievable TOF localisation}
\acr{NECR}{\textbf{noise-equivalent count rate}---a count rate corrected for the noise contributed by randoms and scatter, the standard scalar measure of a scanner's sensitivity}
\acr{FWHM}{\textbf{full width at half maximum}---the conventional measure of spatial, energy and timing resolutions}
\acr{FWTM}{\textbf{full width at tenth maximum}---a similar measure of spatial, energy and timing resolution as FWHM}
\acr{SUV}{\textbf{standardised uptake value}---a semi-quantitative measure of tracer uptake}
\acr{NEMA}{\textbf{National Electrical Manufacturers Association}---publisher of the NU 2 standard by which PET scanner performance is measured}
\acr{DOI}{\textbf{depth of interaction}---the radial depth at which a photon interacts within a crystal, unresolved in conventional designs and a source of parallax error}
\acr{FOV}{\textbf{field of view}---the region imaged by the scanner}
\acr{AFOV}{\textbf{axial field of view}---the extent of the FOV along the patient axis}
\acr{LAFOV}{\textbf{long axial field of view}---scanner designs extending the AFOV towards total-body coverage}
\acr{UBP}{\textbf{unfiltered backprojection}---primitive image reconstruction method which simply backprojects all LORs through a voxelwise volume, produces a distinctly blurred (by  $1/r$) image}
\acr{FBP}{\textbf{filtered backprojection}---UBP preceded by a ramp filter that removes the $1/r$ blur, a standard simple analytical reconstruction algorithm}
\acr{MLEM}{\textbf{maximum-likelihood expectation maximisation}---an iterative reconstruction algorithm}
\acr{OSEM}{\textbf{ordered-subset expectation maximisation}---an accelerated variant of MLEM, standard in clinical practice}
\acr{PSF}{\textbf{point spread function}---the imaged response to a point source, modelled within reconstruction to recover resolution}
\acr{ROI}{\textbf{region of interest}}
\acr{CRLB}{\textbf{Cramér--Rao lower bound}---an information-theoretic bound on the minimum variance achievable by any unbiased estimator of a deterministic parameter}
\acr{MC}{\textbf{Monte Carlo}---stochastic simulation of a process involving random variables; here, of photon transport and detection}
\acr{GATE}{\textbf{Geant4 Application for Tomographic Emission}---a Monte Carlo simulation toolkit for emission imaging}
\acr{SSS}{\textbf{single scatter simulation}---a simulation method for estimating the expected scatter in an acquired PET scan, informed by an attenuation map}
\acr{SPECT}{\textbf{single-photon emission computed tomography}---the single-photon counterpart to PET, in which tomographic reconstruction was first demonstrated}
\acr{PSMA}{\textbf{prostate-specific membrane antigen}---the target of a widely used class of radioligands}
\acr{PSNR}{\textbf{peak signal-to-noise ratio}}

% Terms with no acronym. Same list, same sort -- see the header comment.
\term{annihilation photons}{the two \keV{511} photons produced when a positron annihilates with an atomic electron.}
\term{positron range}{the displacement between the point of positron emission and the point of annihilation.}


\chapter*{Acronyms and Terms}
\addcontentsline{toc}{chapter}{Acronyms and Terms}
% xltabular = longtable (breaks across pages) + tabularx's X column, needed
% once the list outgrew one page. Column 1 is a fixed p, not an l: a multi-word
% \term would stretch an l column and crush the meaning beside it.
\begin{xltabular}{\textwidth}{@{}>{\raggedright\arraybackslash}p{0.23\textwidth} @{\hspace{1.5em}} X@{}}
\toprule
\textbf{Term} & \textbf{Meaning}\\
\midrule
\endfirsthead
\toprule
\textbf{Term} & \textbf{Meaning}\\
\midrule
\endhead
\bottomrule
\endlastfoot
\printacronyms\\
\end{xltabular}
